\PassOptionsToPackage{svgnames}{xcolor}
\documentclass[11pt,a4paper,reqno]{article}
\usepackage{amssymb,amsmath,amsthm}
\usepackage[pdftex]{graphicx}
\usepackage{wrapfig,subfig}
\usepackage[margin=1.4in]{geometry}
\usepackage{xcolor}
\usepackage{tikz}

\usepackage{tcolorbox}
\usepackage{lipsum}
\tcbuselibrary{skins,breakable}

\usepackage{tikz}
\usetikzlibrary{shapes,arrows}

\usepackage{algpseudocode}
\usepackage{algorithm}

\usepackage{dsfont}
\numberwithin{equation}{section}
\usetikzlibrary{intersections}
\usetikzlibrary{calc}

\newcommand{\x}{{\mathbf x}}
\newcommand{\y}{{\mathbf y}}
\newcommand{\etab}{{\mathbf \eta}}

\newtheorem{thm}{Theorem} 
\newcommand{\Cr}{\textcolor{red}}

\newenvironment{myblock}[1]{%
    \tcolorbox[beamer,%
    noparskip,breakable,
    colback=LightBlue,colframe=DarkBlue,%
    colbacklower=DarkBlue!75!LightBlue,%
    title=#1]}%
    {\endtcolorbox}



\usepackage{listings}
\usepackage[square]{natbib}

\setlength\parindent{0pt}

\begin{document}

\bibliographystyle{plainnat}

\title{Early Warning Signals for Tipping Points in Geophysical Systems\\ {\small Scientific Report for Monitoring Committee Metting \# 4}}
\author{Joshua Prettyman\\ University of Reading}
\maketitle

% Define block styles
\tikzstyle{decision} = [diamond, draw, fill=blue!20, 
    text width=4.5em, text badly centered, node distance=3cm, inner sep=0pt]
\tikzstyle{block} = [rectangle, draw, fill=blue!20, 
    text width=10em, text centered, rounded corners, minimum height=4em]
\tikzstyle{line} = [draw, -latex']
\tikzstyle{cloud} = [draw, ellipse,fill=red!20, node distance=3cm,
    minimum height=2em]
    
\section*{Overview}

The project intends to create a multi-purpose tool to predict a tipping point (critical transition) in a dynamical system, specifically a geophysical system, given time series data as input, both in one-dimensional systems and in higher dimensions.\\

Table of contents:
\begin{itemize}
\item[\ref{sec_tippingpoints}.] Introduction to tipping points
\item[\ref{sec_EWS_methods}.] Analytical derivation of one-dimensional Early Warning Signals (EWS) 
\item[\ref{sec_gridded}.] Extension to gridded and multi-variate data (EOF/other)
\item[\ref{sec_hurricanes}.] Application to tropical cyclones
\end{itemize}

\section{Introduction to tipping points}
\label{sec_tippingpoints}

This project is concerned with tipping points, or critical transitions \citep{lenton2008, scheffer2009}, in dynamical systems, and methods that enable us to predict these. \cite{livina2011} identifies three types of system changes which may be the cause of a tipping event:
\begin{enumerate}
\item A forced transition, where the potential function stays the same shape, but is moved, so that there is a trend (or possibly a sudden jump) in the time series.
\item A noise-induced transition, where the system shifts from one stable state to another due to a large initial perturbation.
\item A genuine bifurcation, where the shape of the potential function is changed. 
\end{enumerate}
Given a set of time series data (e.g. a measurement variable of a physical system or computational model) where a transition or bifurcation has occurred, we wish to detect an Early Warning Signal (EWS), which should be measurable before the tipping point. This EWS can then be used to predict future tipping points.\\

\cite{livina2011} study a system such as
\begin{align}
\begin{split}
\label{noise_wells}
\dot{z}(t) & = -(z^4 - 2z^2)' + \sigma(t)\eta \\
\sigma(t) & = -0.00007t+1.50045
\end{split}
\end{align}
where a double-well-potential system effectively becomes a single-well system as the noise term becomes so small that it is practically impossible for the system to change state. In section \ref{sec_pitchfork} we study a system with a genuine `pitchfork' bifurcation, we also note the test of EWS on artificial data made from the concatenation of several short time series with different noise structures \citep{livina2012}. Thus, there are many possible time series on which the EWS may be tested.

\section{Analytical derivation of one-dimensional EWS}
\label{sec_EWS_methods}

\cite{dakos2008} lists several EWS techniques which have been attempted: analysis of spectral properties \citep{kleinen2003}, degenerate fingerprinting \citep{held2004} and a modification of Detrended Fluctuation Analysis (DFA) \citep{livina2007}.\\

A key concept involved in these methods is that of resilience or recovery rates \citep{scheffer2001, dakos2012}, that is, how quickly the system returns to the stable state having received a small perturbation. The theory states that the recovery rate will slow down as the system approaches a bifurcation point \citep{scheffer2009}, known as critical slowing down. As an example we take the system in equation \ref{fork_system}:
\begin{align}
\dot{z}(t) & = -\frac{\partial}{\partial z}(z^4 + \left(3-t\right)z^2) + \eta_t
\label{fork_system}
\end{align}
where $\eta_t$ are independent and Gaussian. This system is trapped within a single well of attraction (at $z=0$) for $t\leq 3$, at $t=3$ a bifurcation occurs creating a double well potential system for $t>3$. The shape of the potential at $t=0$, $t=2$ and $t=3$ is shown in figure \ref{three_shapes}. Intuitively, the system will take longer to return to the $z=0$ equilibrium, following a small perturbation, as the sides of the potential well become less steep. The hypothesis is that this critical slowing down will be seen in a wide variety of dynamical systems approaching a bifurcation.\\

\begin{figure}[H]
\centering
\includegraphics[width = \textwidth]{../../notes/images/three_potential_function_shapes.png}
\caption{Showing the shape of the system potential of the system in equation \ref{fork_system} at $t = 0,2,3$. The sides of the potential well become less steep as the system approaches the bifurcation.}
\label{three_shapes}
\end{figure}

\cite{vannes2007} test this hypothesis with a number of ecological models and find that the intensity of the critical slowing down effect is linearly, or almost linearly, related to the distance from the tipping point in all cases.\\

\subsection{Degenerate fingerprinting: autocorrelation}
\label{sec_ACF}
The use of autocorrelation as an EWS is justified by modelling a dynamical system using a one-dimensional autoregressive system:
\begin{align}
y_{n+1} = e^{\kappa\Delta t}y_{n} + \eta_n
\end{align}
\citep{held2004, scheffer2009}. Where $\eta_n$ are independent and Gaussian, the state $y_n=0$ is the equilibrium. The system returns to equilibrium exponentially with rate $\kappa$ (the decay rate). If we simulate critical slowing down by decreasing $\kappa$, the autocorrelation coefficient $\alpha \equiv e^{\kappa\Delta t}$ increases, $\alpha\xrightarrow[\kappa\rightarrow 0]{}1$. It is also shown that as the autocorrelation increases so does the variance, thus detecting an increase in variance provides another EWS. \\

\cite{held2004} use this observation to argue that in a multi-dimensional system one should project onto the first EOF (the basis vector for which the variance of the system is maximised) as this is the one-dimensional basis in which the rise in autocorrelation (and also variance) will occur. We note that the basis in which the variance is maximimal may not be the same as the basis in which the variance is increasing. At least, this is not immediately obvious a priori and may be wanting further investigation. An investigation of this is an ongoing project, progress is documented in section \ref{sec_gridded}. In some dynamical systems it may be that while autocorrelation rises before a tipping point, variance does not \citep{gsell2016}, this property may also provide fruitful investigation.\\

The autocorrelation of a one-dimensional time series $X=\{x_k\}_{k=1}^{N}$ is estimated using the autocorrelation function (ACF) formula:
\begin{align}
\label{eqn_ACF_formula}
\text{ACF}_{l}(X) = \frac{1}{(N-l)\sigma^2}\sum_{j=1}^{N-l}(x_{j}-\mu)(x_{j+l} - \mu)
\end{align} 
\citep{held2004}, where $l$ is the lag and $\mu$ and $\sigma$ are the mean and standard deviation of the series, it is likely in applications that $\mu$ and $\sigma$ are also estimated from the data, using the standard formulae. Other measures of autocorrelation such as the Mann-Kendall coefficient may be used \citep{yue2002}.\\

The rise in ACF --the ``fingerprint'' of the bifurcation-- is  detected by calculating the ACF of a segment of the time series, say $\{X_i\}_{i=m-k}^{m}$, then sliding this window by one point ($m-k+1$ to $m+1$) and calculating the ACF of this new time series segment. 


\subsection{DFA and spectral properties}

Detrended Fluctuation Analysis \citep{kantelhardt2001} is used by \cite{livina2007} as an EWS, similarly to autocorrelation. Taking time series data $X$ of length $N$, the DFA method involves calculating the cumulative sum of $X$, splitting this into $\lfloor N/s \rfloor$ non-overlapping segments and in each segment `detrending' by subtracting an order-$n$ polynomial fit. These detrended segments are named $\{Y_i\}$. The formula in equation \ref{DFA} gives the $s$-dependent fluctuation coefficient $F^{(n)}(s)$:
\begin{align}
F^{(n)}(s) = \sqrt{ \frac{1}{\lfloor N/s \rfloor} \sum_{i=1}^{\lfloor N/s \rfloor} \text{Var}(Y_i)}.
\label{DFA}
\end{align}
Using the relation $F^{(n)}(s) \propto s^{\alpha_n}$, we can perform the algorithm with many values of $s$ to arrive confidently at a value for $\alpha_n$. This value, which we will call $\text{DFA}_n(X)$, is used as an EWS in the same way as the ACF coefficient.\\

Spectral properties of the time series may also be used as indicators \citep{kleinen2003}, in particular the scaling exponent of the power spectrum.  The use of spectral properties as an EWS will be useful when the tipping point is associated with a change in the structure of the noise, or the stationary system becoming non-stationary. However, a shift from a dominance of short-scale memory to long-scale memory, seen in the measurement of the power spectrum scaling exponent (PSE), will be associated with a rise in autocorrelation, and so in this respect the PSE and the ACF are related. 

\subsection{Comparison of one-dimensional EWS methods}
\cite{livina2012} compare the performance ACF and DFA degenerate fingerprinting methods on artificial  data. The potential analysis method \citep{livina2011} is also used for comparison. The ACF and DFA methods are tested of an artificial time series produced by concatenating several series of increasingly correlated noise, from pure white noise to Brownian motion. This is an example of a time series which experiences a rise in autocorrelation simultaneously with a decrease in variance. 
\subsubsection{Model with a pitchfork bifurcation}
\label{sec_pitchfork}
Here we study the pitchfork bifurcation seen in the model of equation \ref{fork_system}. The equation used here is on a different time scale:
\begin{align}
\dot{z}(t) & = -\frac{\partial}{\partial z}\left(z^4 + \left(3-\frac{t}{200}\right)z^2\right) + \eta_t
\label{pitchfork}
\end{align}
so that the bifurcation occurs at $t=600$. Figure \ref{pitchfork_plot} shows 100 instances of this model plotted together. The noise is sufficiently large that the two branches are not distinguishable until at least $t=650$. 
\begin{figure}
\includegraphics[width = \textwidth]{figures/one_and_many_timeseries.png}
\caption{\label{pitchfork_plot} 100 instances of the model in equation \ref{pitchfork}.}
\end{figure}
In figure \ref{ACF_all} we plot the ACF$_1$ statistic (calculated using a 200-point sliding window) for the 100 instances. We can clearly see a increasing trend in the mean (dashed line), visible even at around $t=400$, but this is not so clear for any individual time series. For example, in model run number 80 (picked out in black) there is an increasing trend very early ($t=250$ to $400$) but this is followed by a decreasing trend, before the expected rise at $t=500$. Figure \ref{comparison} shows the mean, for all 100 time series, of the ACF, DFA and PSE indicators. All show the expected rising trend before the bifurcation, but we must note that in measurements of geophysical systems it is not always possible to run 100 experiments in order to get a useful EWS.  
\begin{figure}
			\includegraphics[width = \textwidth]{figures/one_and_many_ACF.png}
			\caption{\label{ACF_all} The ACF lag-1 estimator is calculated for 100 model runs upto the point of bifurcation ($t=600$). }
			\end{figure}				

			\begin{figure}
			\includegraphics[width = \textwidth]{figures/method_comparison.png}
			\caption{\label{comparison} A comparison of the three autocorrelation indicators discussed. In each case the indicators are calculated for all 100 series, using a widow of 200 points, and then the mean of all is taken as in Figure \ref{ACF_all}. The scaling exponent and DFA coefficient have been scaled linearly to fit on the same axis.}
			\end{figure}	

\subsubsection{Application to real data}
\label{sec_realdata}

In section \ref{sec_hurricanes} we apply the one-dimensional EWS methods to sea-level pressure data time series, it would also be good to apply these methods to some other, similar data for comparison. \cite{dakos2012} compare the results of applying an ACF method to various geophysical time series of where dramatic shifts can be seen. These methods are also used in biological and palaeoclimate data \citep{scheffer2001, livina2011}. 


\section{Extension to multivariate and gridded data}
\label{sec_gridded}

The methods described so far (degenerate fingerprinting, potential analysis) are applied only to a single one-dimensional time series. Since this project aims to study geophysical variables we must consider cases where we have more than one time series, either because:
\begin{enumerate}
\item There is more than one measured variable that may affect an EWS.
\item A variable is measured at several locations on a grid. 
\item Both of the above simultaneously. 
\end{enumerate}
In the first case it may suffice to attempt an EWS detection on each time series individually but one can envision a scenario where it is not possible to detect an EWS in $X(t)$ nor $Y(t)$ but, considered together in some way, an EWS is visible. It may therefore be worthwhile to investigate the possibility of a 2D (or multidimensional) analogue to the one-dimensional methods described in section \ref{sec_EWS_methods}. In the case of gridded data it is common to use Empirical Orthogonal Functions (EOF)\footnote{also know as Principal Component Analysis (PCA)} to reduce the dimensionality \citep[Chapter~13]{vonstorch}, but it may be worth investigating other methods. It is also possible to use the EOF method for multivariate data.

\subsection{A 2D analogue of the ACF fingerprinting method}
\label{sec_multivariate}
As in section \ref{sec_ACF}, we may linearly approximate many dynamical systems using an autoregressive model:
\begin{align}
x_{t+1} = ax_t + c + \varepsilon_t,
\end{align}
and we can estimate the lag-1 autocorrelation coefficient $a$ using equation \ref{eqn_ACF_formula}, or 
\begin{align}
a = \frac{E(x_{t+1}x_t)-\bar{x}^2}{E(x_{t}^2)-\bar{x}^2}
\end{align}
It is proposed that the same method could be applied to a higher dimension system
\begin{align}
\mathbf{x}_{t+1} = A\mathbf{x}_t + \mathbf{c} + \mathbf{\varepsilon}_t
\end{align}
using a straight-forward translation of the one-dimensional formula:
\begin{align}
A= \frac{E(\x_{t+1}\x_t^\top)-\bar{\x}\bar{\x}^\top}{E(\x_t\x_t^\top)-\bar{\x}\bar{\x}^\top}
\end{align}
\citep{williamson2015}, where the fraction is an obvious abuse of notation. For a system of three variables, the matrix $A$ has nine elements, thus we have not done much to reduce the dimension of the problem. It is proposed in the same paper that the eigenvalues of the Jacobian of $A$ should be studied. This method is applied to the system in equation \ref{eqn_homoclinic_system}
\begin{align}
\begin{split}
\dot{x} &= y + \varepsilon^{(x)}\\
\dot{y} &= \mu - x^2 + \varepsilon^{(y)}\\
\mu &= 0.5-0.005t
\end{split}
\label{eqn_homoclinic_system}
\end{align}
with a homoclinic bifurcation at $\mu(t)=0$ ($t=100$). The system is integrated from $t=0$ to $t=100$ with $\varepsilon$ white noise with standard deviation $\sigma = 0.5$. The solution is given at intervals of $\Delta t=0.5$. The result is shown in figure \ref{ourownonorbit0p5} and the eigenvalue analysis is shown in figure \ref{Jeigenvals0p5}. 
\begin{figure}
\centering
\includegraphics[width = 5in]{../../notes/images/ourownorbit0p5.png}
\caption{Attempt to replicate the results in \citep{williamson2015}, with the noise std of 0.5.}
\label{ourownonorbit0p5}
\end{figure}
\begin{figure}
\centering
\includegraphics[width = 5.5in]{../../notes/images/Jeigenvals0p5.png}
\caption{The Jacobian eigenvalue (both real and imaginary parts) plotted in a sliding window. These are derived from the time series shown in figure \ref{ourownonorbit0p5}}
\label{Jeigenvals0p5}
\end{figure}
The decreasing imaginary part of the principal eigenvalue predicates the bifurcation, as predicted by the analysis, but it is not clear that this could serve as an EWS without working out beforehand which changes to look for, already having perfect knowledge of the equations describing the system. Also, for an $n$-dimensional system we have $2n$ variables to analyse ($n$ eigenvalues, each with real and imaginary parts). However, studying the principal eigenvalue, rather than the $n$ individual components of the system, does at least incorporate part of the information of each variable into a single statistic, even if much information is lost by ignoring the other eigenvalues. We must investigate in which cases the principal eigenvalue provides a useful EWS.

\subsection{Investigation of the EOF method}
\label{sec_EOF}
The use of EOF to reduce dimensionality is common in climate research \citep[Chapter~13]{vonstorch}. The method aims to ``capture'' most of the ``interesting'' behaviour of the system \citep{held2004}. We project the multi-dimensional system onto an orthogonal basis such that the first component of the system in the new basis has maximal variance, and the $n^{\text{th}}$ component has maximal variance given the first $n-1$ basis vectors.  One may then study only the first $k$ components of the system in the new basis, where $k$ is chosen so that, say, 95\% of the total variance is captured. In many applications only the one-dimensional system of the first component is considered.\\

We wish to ask the question: how useful or relevant is the first EOF in the detection or prediction of tipping points? That is, are there cases where the information relevant to tipping, such increasing autocorrelation \citep{scheffer2009}, is not captured by the first EOF? We attempt to answer this question by studying the stochastic system eqn.\ref{sytem_eqn_1},
\begin{align}
\x_{k+1} = B\x_k + S\eta_k,
\label{sytem_eqn_1}
\end{align}
where the $\etab_k$ are column vectors with each element $\sim N(0,1)$, iid. Suppose that $B$ is diagonalisable with real eigenvalues $\mu_{i}$. In the case that all $|\mu_{i}|<1$, $\x_n$ will tend to zero but it will tend to zero more slowly in the direction $\mathbf{v}_{\max}$, the eigenvector corresponding to $\mu_{\max} := \max_{i}|\mu_{i}|$. If we allow $\mu_{\max}$ to vary a critical transition occurs at $|\mu_{\max}|=1$ and $|\mu_{\max}|>1$ will cause the $\x_n$ to tend to infinity in the direction $\mathbf{v}_{\max}$. If we project the series $\x_n$ onto the vector $\mathbf{v}_{\max}$, we expect to see an increase in autocorrelation as $\mu_{\max} \rightarrow 1$. If the first EOF is to capture this autocorrelation and provide a good predictor of the tipping point, we would expect that the first EOF score, $\mathbf{w}_1$, is very close to $\mathbf{v}_{\max}$. \\

\subsubsection{An empirical calculation of EOF}
We attempt to calculate empirically the eigenvectors of the covariance matrix of the series $\{\x\}$ described by equation \ref{sytem_eqn_1}. In the first instance we assume that $B$ is constant with eigenvalues $\mu_i<1$ $\forall i$ so that there is no tipping point and the long-term mean of $\{\x\}$ is zero. We also assume that $B$ is diagonal, which is equivalent --via a change of basis-- to saying that $B$ is diagonalisable. \\

We first write out the terms of equation \ref{sytem_eqn_1} extrinsically in terms of $\x_0$:
\begin{align}
\x_k &= B^k\x_0 + \sum_{i=1}^{k} B^{k-i}S\etab_i \\
&= \sum_{s=0}^{\infty} B^{s} S \etab_{k-s} - B^k \left(\sum_{s=0}^{\infty} B^{s} S \etab_{-s} - \x_0\right).
\end{align}
And we replace $k$ with $k+N_0$ where $N_0$ is large enough so that $B^{N_0} \approx 0$ and the mean of $\{x_i\}_{i=0}^{N_0}$ is also approximately zero. It is then possible to calculate the expectation of the covariance of the time series:
\begin{align}
\mathbb{E}(C) &= \frac{1}{N}\sum_{k=0}^{N-1}\left[\sum_{s=0}^\infty \sum_{r=0}^\infty B^s S\mathbb{E}\left[(\eta_{k+N_0-s})(\eta_{k+N_0-r})^\top\right] S^\top (B^\top)^r\right]\\
& = \sum_{s=0}^\infty B^s S S^\top (B^\top)^s \\
D :=\mathbb{E}(C) &= SS^\top + B D B^\top
\end{align}
It is then possible to calculate the eigenvalues and eigenvectors of $D$. We here assume that $B$ is diagonal and that the system is two-dimensional, with $b_{11}$ very close to $1$ and greater than $b_{22}$. We note that 
\begin{align}
D = \left[ 
\begin{array}{cc}
\frac{s_{12}^2+s_{11}^2}{1-b_{11}^2} & \frac{s_{12}s_{22}+s_{11}s_{21}}{1-b_{11}b_{22}}\\
\frac{s_{12}s_{22}+s_{11}s_{21}}{1-b_{11}b_{22}} & \frac{s_{22}^2+s_{21}^2}{1-b_{22}^2}
\end{array}
\right].
\label{D_for_B_diag}
\end{align}
In general, the eigenvalues of symmetric matrix 
\(
\left(
\begin{array}{cc}
a&b\\
b&c
\end{array}
\right)
\)
are 
\begin{align}
\frac{1}{2}\left(a+c\pm\sqrt{(a-c)^2+4b^2}\right)
\label{general_eigenvalues}
\end{align}
and since everything is positive the positive square root will give the largest eigenvalue. Note that for $b=0$ the difference between the two eigenvalues is $|a-c|$ but for non zero $b$ this difference increases -- the eigenvalues become more distinct. For matrix $D$ this $b=0$ corresponds to the situation where $S$ is diagonal or reverse-diagonal, or one of the rows or columns is zero. The eigenvector corresponding to this largest eigenvalue (the principal eigenvector) is
\begin{align}
\left(
\begin{array}{c}
a-c + \sqrt{(a-c)^2+4b^2}\\
2b
\end{array}
\right).
\label{largest_eigenvector}
\end{align}

We substitute the terms from equation \ref{D_for_B_diag} into equation \ref{largest_eigenvector} and expand out the roots and reciprocals upto terms $\mathcal{O}((1-b_{11})^2)$, since $(1-b_{11})\approx 0$. We find that the first component is
\begin{align}
 v_1 = \frac{2p}{(1-b_{11}^2)} + 
\frac{-2q}{(1-b_{22}^2)}+
(1-b_{11}^2) \left[ \frac{2r^2(1+ b_{22})^2 -  q^2 }{p(1-b_{22}^2)^2}\right]
\end{align} 
and the second vector component is
\begin{align}
v_2 = 2r\left[\frac{1}{1-b_{22}} -
\frac{b_{22}(1-b_{11})}{(1-b_{22})^2}
\right]
\end{align}
where $p := s_{11}^2+s_{12}^2$, $q := s_{21}^2+s_{22}^2$ and $r := s_{12}s_{22}+s_{11}s_{21}$. We expect this eigenvector, $\mathbf{v} = (v_1, v_2)^\top$, to be approximately $(1,0)^\top$, since the ``interesting'' behaviour is happening in this (the $x$) direction. The system tends to zero much more quickly in the $y$ direction. If $S$ is diagonal ($r=0$) then this is always the case, otherwise it may be the case that this will be closer to the direction of $(0,1)^\top$, but this is only if $s_{11},s_{12}$ are very small and $s_{21},s_{22}$ are very large. In this case it is already very obvious that the largest variance will be in the $(0,1)^\top$ projection, it is simply a case of a very large noise obscuring the signal (in this case the `signal' is the impending tipping point as $b_{11}$ `approaches' 1).

\subsubsection{Further comments}

So far in this section we have shown that the EOF assumption is valid in the case of the system of equation \ref{sytem_eqn_1}, subject to several assumptions. The next step is to consider the cases of a three (or higher) dimensional system, an $n$-dimensional system, where $B$ is non-diagonalisable, and where the system is non-stationary. That is, where $B$ is has a time dependence and one of its eigenvalues approaches 1 (from below). \\

It will probably not be possible to calculate the EOF explicitly where $B$ is non-diagonalisable, and when $B$ is time-dependent, therefore we may infer patterns from repeated simulations in Matlab.

\section{Application to Tropical Cyclones}
\label{sec_hurricanes}
Tropical cyclones are a good example of something happening. Looking at the sea-level pressure in one location, all is fairly constant, then suddenly there is a drop. One would expect that there might be a build up of correlations before the event. It is possibly not unreasonable to expect that an indication that the event is approaching is detectable before the event is detected. That is, that the outcome of an analysis of the data will reveal something such as ``indicator $\beta$ becoming greater than $0.95$'' before it is possible to say ``the pressure has decreased more than would be expected by random chance, indicating that a storm system is close by.''\\

The analogy to a double well potential two-state system is not at all obvious here, as it is in the hot-Earth/ice-age example in \cite{livina2011}. It fact this may not provide any analogy at all. However, the time series look superficially similar, and this is also true of many economic and physical systems where sudden catastrophe occurs.\\

\subsection{Choice of data}

First, we have the option of where to look (geographically). We also have the option, which variable to use: Sea-Level Pressure (slp), wind speed, etc., if we wish to use gridded data, or more than one variable, we will need to consider the methods of section \ref{sec_gridded}. For this investigation, the HadISD dataset\footnote{Raw weather station data. Available online from the Met Office.} was used. Several tropical cyclones from the past 50 years were selected, using a list of strongest (category 4 or 5) storms, and the data was taken from the weather station closest to the point where the storm made landfall.\\

It was discovered that all slp data contains oscillations of roughly 12-hour wavelength. We experimented removing these oscillations using a simple de-seasonalising method:
\begin{align}
x_i \longrightarrow x_i - \text{mean}\left( [x_{i+12k}]_{k = 0, \pm 1, ...} \right)
\end{align}
and also using a Fourier analysis approach to remove peaks from the power spectrum. These methods had almost identical results, but the Fourier analysis is obviously more robust to changes in the frequency of the oscillations. This method was used to remove these oscillations from all the HadISD data.\\ 

We are also faced with the problem that all cyclones are very different. Probably the most useful way to classify the cyclones is by whether they are building, or past their peak, at the moment they land on the coast.


\subsection{Results of degenerate fingerprinting methods}

Some experiments were made using the eigenvector method \citep{williamson2015}, but these are not yet well-understood. A better understanding of the method, and the scenarios in which the method will work, is required before conclusions can be drawn. Some experiments were also made using gridded (reanalysis) data where the basic EOF method was applied to yield a one-dimensional time series. Here we show results from the use of the HadISD dataset, for consistency. \\

The ACF lag-1 of the sea-level pressure is calculated in a sliding window of 120-240 points. In some cases we see an expected rise in the ACF just before the cyclone event, however in other cases the value does not rise, see figure \ref{Charlie_and_Andrew}, DFA fingerprinting is similarly inconclusive. However, the autocorrelation does appear to be worthy of further study. The sliding ACF was calculated for many cyclones and for many windows of data where no cyclone occurred. In the cyclone time series, the ACF had much larger variances and lower means than in the control. Therefore, the ACF does in some way give an indication of whether a cyclone may be approaching, although these results may simply be correlated to other weather phenomena exclusive to the summer months (when cyclones occur). \\

The PSE (power spectrum exponent) method appears to give a somewhat better EWS. Figure \ref{fig_PSE_allstorms}  shows slp data from 20 cyclones and the PSE calculated in a 50 point (50 hour) sliding window. There is a clear rising trend in the mean PSE, and this trend can also be seen in the PSE for most (but not all) of the individual time series.  



\begin{figure}
\centering
\includegraphics[width = \textwidth]{../../notes/images/CharlieandAndrew.png}
\caption{slp data for hurricanes Charlie and Andrew (left) and the ACF of this data measured in a 240-point sliding window. Contrast the rising ACF for Charlie with the steep decrease for Andrew.}
\label{Charlie_and_Andrew}
\end{figure}

\begin{figure}
\centering
\includegraphics[width = \textwidth]{../../notes/images/SEX_all20storms.png}
\caption{\label{fig_PSE_allstorms} PSE calculated in a 50 point sliding window. The black (dashed) line shows the mean of all PSE values.}
\end{figure}

\subsection{A stochastic model}
There is a need to build a stochastic model of the approaching-cyclone time series. We suggest that the model be a double-well potential with added daily and seasonal oscillations and added noise, parametrised using the Unscented Kalman Filter as described by \cite{kwasniok2009}.

\subsection{The Null Hypothesis}
Our work so far has been concerned mainly with identifying EWS retrospectively, but we have in mind that an EWS should be used predictively. We must decide which is the null hypothesis that our method attempts to answer:
\begin{enumerate}
\item Will there be a cyclone?
\item Will there be a cyclone in the next [three days]?
\item Given that we know there will be a cyclone in the next [ten days], will there be a cyclone in the next [three days]?
\item Given that we know there will be a cyclone in the next [ten days], when will the event happen?
\item In which direction the cyclone is likely to propagate, given a number of observational nodes?
\end{enumerate}
It will be very difficult to answer most of these questions. 

\section{Future work}

At the end of section \ref{sec_tippingpoints} we note that there are many possible time series on which the EWS may be tested. Classifying the cases in which the various EWS may be used is a possible future aim of this project, although it seems unlikely that much will be achieved in this direction. We note that it is an even more difficult proposition when we also consider time series from real-life measurements of geophysical variables.\\

In section \ref{sec_EWS_methods} we note that \cite{vannes2007} test the hypothesis that critical slowing down is a good EWS in a variety of ecological models. It may be necessary to test this further ourselves.\\

We note that the estimator for the auto correlation coefficient in equation \ref{eqn_ACF_formula} may not be the best estimator to use. An approach using a maximum likelihood estimator has been discussed and is a possible future project.\\

In section \ref{sec_gridded} we motivate the development of a multi-variate EWS method with the hypothesis that there exist systems where it is not possible to detect an EWS in $X(t)$ not $Y(t)$ but, considered together in some way, an EWS is visible. We should construct such a system as a test.\\

We also plan to continue the study of tropical cyclones, as noted in the text of section \ref{sec_hurricanes}, and to further pursue the study of the EOF method.

\newpage
\bibliography{../../notes/bibliography.bib}

\end{document}
